\chapter{Design of a Frequent Batch Auction Engine}
\label{ch:design}

This chapter breaks a batch auction engine into five design blocks, lists the
options available in each block, and explains from theory how participants
should be expected to behave under each option. The most consequential of
these behaviours, the boundary repricing race that a predictable close
produces, is developed in full as part of Block~1
(\Cref{sec:priceimprovement}), and \Cref{sec:consequences} closes the
chapter with a summary of every option and its consequences. The chapter
uses the notation of \Cref{ch:background}, in particular the clearing
mechanism defined in \Cref{sec:fbanotation}.

\section{The Engine as a Set of Design Blocks}
\label{sec:blocks}

A matching engine takes a stream of timestamped messages and produces trades.
A continuous engine handles each message completely before reading the next.
A batch engine first collects messages for a while, then processes the whole
set at once. This one change forces five decisions that a continuous engine
never has to make, because serial processing settles them automatically:

\begin{enumerate}
  \item how long the collection period $\tau$ lasts, and whether its end can be
  predicted;
  \item what participants can see while orders are being collected;
  \item whether all orders compete in one pool;
  \item how the trading price is computed and how the orders on the longer
  side are served;
  \item how the trades are settled and what happens to orders that did not
  trade.
\end{enumerate}

\Cref{fig:pipeline} shows how these blocks fit together in one batch cycle.

\begin{figure}[htbp]
\centering
\fbox{\begin{minipage}{0.92\textwidth}
\vspace{0.4em}
\centering
\ttfamily\small
\begin{tabular}{c}
\textbf{[1] Accumulation phase, duration $\tau$}\\
orders, cancels and modifies buffered; nothing executes\\
$\downarrow$ \quad \emph{governed by Block 1 (interval) and Block 2 (disclosure)}\\[0.5em]
\textbf{[2] Gate closes}\\
buffer frozen; late messages routed to the next batch\\
$\downarrow$\\[0.5em]
\textbf{[3] Partition into auctions}\\
one pool, or maker/taker and retail/institutional partitions\\
$\downarrow$ \quad \emph{governed by Block 3 (segregation)}\\[0.5em]
\textbf{[4] Clearing}\\
solve for $\clr$ maximising executed volume\\
$\downarrow$ \quad \emph{governed by Block 4 (clearing)}\\[0.5em]
\textbf{[5] Rationing}\\
allocate the long side among eligible orders\\
$\downarrow$ \quad \emph{governed by Block 4 (rationing)}\\[0.5em]
\textbf{[6] Settlement and residual handling}\\
transfers executed; unmatched volume rolled over or cancelled\\
$\downarrow$ \quad \emph{governed by Block 5}\\[0.5em]
\textbf{[7] Publication} $\rightarrow$ next batch begins
\end{tabular}
\vspace{0.4em}
\end{minipage}}
\caption[The batch cycle and its design blocks]{One batch cycle of a
frequent batch auction engine, annotated with the design block that
governs each stage.}
\label{fig:pipeline}
\end{figure}

\section{Block 1: The Batch Interval}
\label{sec:interval}

\subsection{The Length of \texorpdfstring{$\tau$}{tau}}

The interval length decides how much of the speed advantage disappears. A
speed difference of $\delta$ between two participants only matters if it
changes which batch their orders land in. If orders arrive at random times,
this happens with probability $\delta/\tau$. With $\delta$ of a few
microseconds --- the time scale on which the races of
\citet{aquilina2022quantifying} are decided --- and $\tau$ around
100~milliseconds, the probability is only $10^{-4}$ to $10^{-5}$. Speed is not
banned; it simply stops paying.

The price is delay. Every participant waits on average $\tau/2$ before its
order can trade, and up to $\tau$ in the worst case. This is harmless for an
investor thinking in minutes, but it hurts a market maker holding inventory,
because the price can move during the interval and the order cannot be
withdrawn in time. The interval should therefore be long enough to remove
speed differences and short enough to keep prices current. The right value
depends on the instrument:

\begin{itemize}
  \item \textbf{Lower bound:} $\tau$ must be clearly longer than the
  differences in network delay between participants. Otherwise faster
  participants still reach the current batch more often and a weaker speed
  race remains.
  \item \textbf{Upper bound:} $\tau$ must be short enough for new information
  to reach prices quickly. \citet{zoican2021speed} show that more frequent
  auctions do not always improve market quality, and that very long intervals
  let orders become stale.
  \item \textbf{Liquidity:} in a thinly traded instrument a short interval may
  collect no matching orders at all, leaving batches empty, so the market
  behaves like a continuous one with an added delay.
\end{itemize}

\subsection{Fixed versus Randomised Closing}

If everyone knows when the batch closes, that moment becomes a target. A fast
participant can wait until the last possible moment, see as much information
as possible, and only then submit, reprice or cancel. This is the
\emph{boundary repricing race} (sometimes called boundary sniping). It is the
same speed race as in continuous trading, moved to the end of the interval;
what is won is a better share of the surplus of a batch whose contents are
now partly known, rather than a single stale quote. \Cref{sec:priceimprovement}
derives it formally; the options below differ in how exposed they are to it.

\begin{itemize}
  \item \textbf{Fixed closing time.} Simple, easy to audit, and compatible with
  clock-aligned reporting, but easy for participants to plan around and fully
  exposed to the boundary repricing race. This is the option the simulated
  engine implements (\Cref{tab:fbachoices}).
  \item \textbf{Random closing time.} The batch closes at a time drawn from a
  publicly announced range, for example uniformly between $\tau_{\min}$ and
  $\tau_{\max}$. Waiting until the last moment now risks missing the batch,
  so it becomes a gamble instead of the best strategy. The drawbacks are that
  participants cannot time their orders precisely and that the random draw
  must be verifiable, which is hard in a decentralised setting.
  \item \textbf{Closing triggered by volume or events.} The batch closes when
  enough volume has accumulated or a market event occurs. This adapts the
  interval to activity, but because the closing depends on what participants
  do, it can be manipulated: a participant can trigger the close by
  submitting volume.
  \item \textbf{Random end after a minimum duration.} The batch runs for a
  fixed minimum time and then ends at a random moment within a short
  window. Xetra's closing auction works this way: after a minimum call
  phase it ends at a random time, precisely to make the price harder to
  manipulate \citep{deutscheboerse2026xetra}. The schedule stays predictable
  to within that short window, while the exact final moment, the one a
  late submitter would target, is not.
\end{itemize}

\subsection{Consequence of the Time Window: Price Improvement and the
Boundary Repricing Race}
\label{sec:priceimprovement}

Choosing a time window, and in particular a window whose end everyone can
predict, has a consequence that the option list above only names. This
subsection derives it. The argument is, to the author's knowledge, not
stated in this form in the existing literature. Its components are
individually familiar: the residual value of low latency near a batch
boundary is established by \citet{schlegel2023transaction}, the tendency of
sophisticated traders to submit late in transparent call auctions is
documented by \citet{daures2026click}, and the strategic consequences of
the clearing-price selection rule in uniform-price auctions are studied by
\citet{burkett2020uniform}. What has not been made is the connection
between them. The frequent batch auction literature treats the price
improvement that uniform clearing confers on inframarginal orders as a
benefit of the design; this subsection argues that the treatment rests on a
category error. Price improvement measured on one side of the book is not
a welfare quantity but a transfer, and the rule that determines its size
is manipulable by the timing of order submission. A batch with a
predictable close therefore does not remove the incentive to compete on
speed: it converts a race for priority into a race for proximity to the
batch boundary, independently of the informational motives usually
invoked to explain late submission.

The argument considers a single batch and uses the objects of
\Cref{sec:fbanotation}: the schedules $D(p)$ and $S(p)$, the
volume-maximising set $\mathcal{P}^{*} = [\underline{p}, \overline{p}]$,
the clearing price $\clr = \sigma(\mathcal{P}^{*})$, the executed
quantities $x_k$ with $\sum_{k \in \mathbb{B}} x_k = \sum_{k \in \mathbb{A}}
x_k = V^{*}$, and the buyer and seller surplus $W^{\mathbb{B}}(p)$ and
$W^{\mathbb{A}}(p)$ of \eqref{eq:surplus}.

\subsubsection{Price Improvement Is a Transfer}

\textbf{Total surplus does not depend on which clearing price is chosen.}
For any $p \in \mathcal{P}^{*}$ with allocation $x$ held fixed,
\begin{equation}
\label{lem:invariance}
  W = W^{\mathbb{B}}(p) + W^{\mathbb{A}}(p)
    = \sum_{k \in \mathbb{B}} \pi_k x_k
      - \sum_{k \in \mathbb{A}} \pi_k x_k ,
\end{equation}
which does not depend on $p$.

Expanding the two sums, the terms in $p$ enter as
$p\left(\sum_{\mathbb{A}} x_k - \sum_{\mathbb{B}} x_k\right)$, which
vanishes by flow conservation. Two consequences follow immediately.
First, $\partial W^{\mathbb{B}} / \partial p = -V^{*}$ and
$\partial W^{\mathbb{A}} / \partial p = +V^{*}$: every unit of
price improvement conferred on one side is subtracted from the other.
Second, a measured increase in buy-side price improvement is not evidence
that the mechanism created value: total gains from trade are fixed by the
allocation, and the clearing price only divides them. The metric must
therefore be read as a distributive statistic, reported for both sides
jointly. If $\mathcal{P}^{*}$ were a singleton, the division would be
pinned down by the submitted limits and there would be nothing to
contest; because it is frequently an interval, especially in thin
batches, the surplus split depends on a rule, and any participant able to
influence the position of $\mathcal{P}^{*}$ or of $\sigma$ within it has an
incentive to do so.

\subsubsection{The Last-Mover Advantage}

Suppose the engine operates under the fully open disclosure regime of
Block~2 (\Cref{tab:disclosure}), so that a participant submitting at time
$t$ observes the schedules formed by all orders with $t_k < t$.

\textbf{Late submission weakly dominates.} Suppose a participant can see
the orders already submitted by others, and hence the residual supply or
demand it faces, and that the selection rule $\sigma$ responds
monotonically to submitted limit prices, so that a participant can shift
the clearing price within $\mathcal{P}^{*}$ by adjusting its own limit.
The participant's best response is then to set its limit so that the
clearing price moves to the edge of $\mathcal{P}^{*}$ most favourable to it,
\begin{equation}
\label{prop:lastmover}
  \clr \to \underline{p} \text{ for a buyer,} \qquad
  \clr \to \overline{p} \text{ for a seller,}
\end{equation}
while still keeping its order executed. Because this manoeuvre requires
knowing the other orders, a participant's surplus can only grow with the
information it has when it submits. Submitting later is therefore never
worse than submitting earlier, and it is strictly better whenever the
extra information changes the participant's optimal limit.

The argument is the standard one for sequential-move bargaining over a
fixed surplus, applied to the division established in
\eqref{lem:invariance}. A seller who observes a resting buy limit of
$\pi_{\mathbb{B}}$ knows that any ask weakly below it will execute, and
therefore submits just below $\pi_{\mathbb{B}}$ rather than at its own
valuation, capturing the interval. Truthfulness
(\Cref{sec:auctiontheory}) is not violated in its own terms, since it is
stated for a participant that cannot move the clearing price; what
disclosure does is make the boundaries of $\mathcal{P}^{*}$ computable, and
hence make deviation actionable for any participant whose order is large
enough to reach them.

This identifies the source of the incentive to wait. When several prices
clear the same maximum volume, the choice among them is not neutral: it
transfers surplus between buyers and sellers. Under a monotone $\sigma$,
the participant who moves last can capture this transfer by pushing the
price to the boundary of $\mathcal{P}^{*}$ in its own favour. The prize
goes to whoever submits closest to the deadline with the fullest view of
the book: a reward for speed and timing rather than for providing
liquidity or information.

Two conditions expose a participant to being the donating side of the
transfer: submitting without observing the accumulated batch, as an
uninformed participant does, and submitting early without repricing. Since
the transfer is zero-sum by \eqref{lem:invariance}, every participant
exposed in this way is matched by one with an incentive to exploit it. An
order resting in an open batch is thus a free option written to every
subsequent submitter in the window. This is the same option
\citet{copeland1983information} describe for a resting limit order in a
continuous book: there, it is written to traders reacting to news; here,
the same option is written to traders reacting to the accumulating batch
itself.

\subsubsection{Cancellation, Repricing and the Deadline}

The natural defence against the last-mover advantage \eqref{prop:lastmover}
is to withdraw or reprice an order that has become adversely positioned.
This is where the speed race re-enters.

\textbf{Latency monotonicity.} Let participant $i$ face message latency
$\ell_i > 0$ and let the batch close deterministically at $\tau$, so that
the window is $[0, \tau)$ as in \Cref{sec:fbanotation}. A message
dispatched at time $t$ is included in the batch if and only if
\begin{equation}
\label{prop:latency}
  t + \ell_i < \tau,
\end{equation}
so the latest observable state on which $i$ can act is that prevailing
just before $\tau - \ell_i$. The information set available to $i$ at the
moment of its final decision is therefore ordered by $\ell_i$, and
expected surplus is non-increasing in $\ell_i$.

The structure of the advantage differs from the continuous case but its
source does not. In a continuous book, lower latency wins a contest for
priority; in a batch, priority within the window has no effect on the
price by construction (and affects the allocation only if the rationing
rule rewards it, \Cref{sec:rationing}), and lower latency instead buys a
later decision point. The fast trader observes more of the batch before
committing, can reprice on more recent information, and can place its
final message closer to the boundary while retaining the same probability
of inclusion. Speed is not neutralised; its payoff is relocated from
ordering to deadline precision. That some speed advantage survives
batching is consistent with the experimental evidence of
\citet{aldrich2020experiments}, who find that investment in speed falls but
does not disappear under batch clearing, and with the model of
\citet{zoican2021speed}.

Three remarks bound the mechanism. Prohibiting cancellation within the
batch does not close the channel: a participant with two-sided access can
still neutralise an adverse position with an offsetting order, a
synthetic cancellation that needs no cancel message at all (true
enforcement would require account-level self-match prevention). Banning
cancellation outright would, if anything, strengthen the last-mover
advantage \eqref{prop:lastmover}, by raising the cost of correcting an
early submission instead of removing it. And the channel is
disclosure-conditional: the informational component of the last-mover
advantage requires an observable accumulating batch and is absent under a
sealed-bid regime.

\subsubsection{Equilibrium Implication and Design Responses}

If late submission weakly dominates and lateness is bounded only by
$\ell_i$, submissions concentrate near the boundary and are ordered by
latency. In the limit the accumulation phase carries no orders and
therefore no information, and the auction degenerates into a sealed tender
conducted in the final microseconds of the window: the batch retains its
length nominally while losing the property the length was chosen to
deliver.

Three design responses follow, one in this block and two in later ones.
Randomising the close time (the random closing option above) makes waiting
costly, since an order dispatched at $t$ is included only if the drawn
close exceeds $t + \ell_i$; the participant then solves an interior
optimum trading the value of information against the risk of exclusion,
which justifies randomisation independently of the sniping argument
usually advanced for it. Restricting disclosure to an aggregate imbalance
or to nothing (Block~2, \Cref{sec:disclosure}) removes the informational
component of the last-mover advantage while leaving bid shading, so the
choice is between two costs rather than between one cost and none.
Finally, and least discussed in the literature, $\sigma$ can be made
insensitive to the marginal submitted limits (Block~4,
\Cref{sec:clearing}), for instance by choosing the point of
$\mathcal{P}^{*}$ closest to an external reference price. This removes the
lever without reducing gains from trade, since by \eqref{lem:invariance}
the choice of price within $\mathcal{P}^{*}$ does not change total surplus
for a given allocation.

The prediction that submissions bunch near the close needs participants
who can see the close and choose when to submit. It cannot be tested in
the replay of this thesis, because the replayed orders were submitted to a
continuous market and never timed against a batch boundary
(\Cref{sec:testable}).

\section{Block 2: Information Disclosure}
\label{sec:disclosure}

The disclosure regime decides what participants can see when they submit.
There are four regimes, from publishing everything to publishing nothing
(\Cref{tab:disclosure}). More secrecy protects submitted orders from being
exploited, but gives participants less price information and encourages
\emph{bid shading}, meaning deliberately submitting a less aggressive price
than one's true valuation.

\begin{table}[htbp]
\centering
\caption[Disclosure regimes and induced behaviour]{Disclosure regimes
during the accumulation phase and their predicted behavioural
consequences.}
\label{tab:disclosure}
\small
\begin{tabularx}{\textwidth}{@{}p{2.9cm} L L@{}}
\toprule
\textbf{Regime} & \textbf{What is published during accumulation} &
\textbf{Predicted behaviour}\\
\midrule
Fully open book & All accumulated orders, continuously updated &
Strong within-batch price discovery; early submitters expose intent and
are exploited by later ones; incentive to submit as late as possible,
which compounds the boundary repricing race (\Cref{sec:priceimprovement})\\
Indicative clearing price & The price at which the batch would clear if
it closed now & Coordination on a common price estimate without exposing
individual orders; some intent leakage through price movement; the
regime of most equity opening and closing auctions\\
Aggregate imbalance & Net excess demand or supply, no price schedule &
Signals direction of pressure and attracts offsetting liquidity without
revealing individual limits; weaker price anchoring\\
Sealed bid & Nothing until the auction clears & No exploitation of
submitted orders and no incentive to delay for informational reasons;
clearing price uncertainty induces bid shading, which can reduce cleared
volume\\
\bottomrule
\end{tabularx}
\end{table}

Disclosure and interval design are linked. An open book gives participants a
reason to submit late, which makes a random closing time more valuable. A
sealed-bid design removes most of that reason and makes a fixed closing time
easier to defend. Secrecy also does not remove information leakage; it only
moves it. Under sealed bidding, the price and volume of one batch become
public before the next one closes, and as $\tau$ gets shorter this sequence of
published prices approaches a continuous price signal.

\section{Block 3: Order Flow Segregation}
\label{sec:segregation}

\subsection{The Unified Pool and Segregation by Participant Class}

In the default design all orders enter one auction. This is the simplest
design to understand and regulate. Its drawback is that a liquidity provider
who cannot tell informed from uninformed traders must quote a price that
reflects the average of the pool, so uninformed participants end up paying for
the losses that informed participants cause. Splitting flow by participant
class, typically retail against institutional, lets the provider quote a
tighter price to the group it believes is uninformed. This is the logic
behind retail wholesaling, the purchase of retail order flow, in equity
markets \citep{easley1996cream}, and it carries over to batches
through one auction per group, or through eligibility rules for each group.
The difficulty is classification. The split is only as good as the label,
labels can be gamed by sending institutional orders through retail channels,
and participants in different groups get different prices for the same asset
at the same moment, which raises fairness and best-execution concerns.

\subsection{Intent-Based Segregation: The Dual Flow Batch Auction}

Jump Trading's Dual Flow Batch Auction (DFBA) \citep{jump2025dfba} splits
orders by \emph{intent} rather than by who submits them. Every order is
labelled either maker flow (supplying liquidity) or taker flow (demanding
it). Each batch then runs two separate auctions, maker buys against taker
sells and maker sells against taker buys, both at one uniform price. Market
makers never trade against each other, which removes trades between two
liquidity providers that add volume but give liquidity to nobody who wanted
it. A fast participant also cannot use taker orders to pick off maker orders
inside the batch, because there is no queue to win.

The design has its own costs. Splitting the pool splits the liquidity, so each
auction clears against a thinner book. Intent labels must be enforced and not
merely declared, otherwise a participant could label itself a maker in one
batch and a taker in the next and collect the benefits of both. And the split
assumes that ``maker'' and ``taker'' mean something, which in a mechanism that
clears everything at the same instant is a rule imposed by the engine, not a
fact about who arrived first.

\subsection{The Maker--Taker Distinction and Fee Design}

This last point leads to a problem the literature on batch auctions has not
solved. In a continuous book the maker is the resting order and the taker is
the order that crosses it, so the distinction depends on arrival order. That
definition is not available when all matched orders arrive in the same
interval and trade at the same time and price. Yet fees, rebates and a large
part of exchange revenue depend on it. This thesis distinguishes four ways
to redefine it, none fully satisfactory; the first follows the DFBA, the
other three are proposed here:

\begin{enumerate}
  \item \textbf{By declared intent}, as in the DFBA. Simple to operate, but
  self-declared unless a separate eligibility rule enforces it.
  \item \textbf{By price aggressiveness.} Orders whose limit price is far from
  the clearing price on the passive side are makers; orders that would have
  crossed a wide range of prices are takers. This can be observed and is hard
  to game without changing one's real exposure, but it penalises participants
  for wanting certainty of execution.
  \item \textbf{By marginality.} The orders that set the clearing price are
  takers; orders priced strictly better than it are makers. This is
  attractive in theory, since the marginal order determines the price, but the
  fee a participant pays then depends on what everyone else does in the
  batch, so it cannot be known in advance.
  \item \textbf{By persistence across batches.} An order that has waited
  through several batches is a maker; a new order is a taker. This rewards
  liquidity that stays available, which is what the rebate is supposed to pay
  for, and it is the only one of the four that measures something lasting
  rather than a single moment.
\end{enumerate}

The choice matters beyond accounting. A rebate is an incentive, so whichever
definition the engine adopts is the behaviour it ends up paying for.

\section{Block 4: Matching, Clearing and Rationing}
\label{sec:clearing}

\subsection{Implementing the Clearing Rule}

Prices can only take values on a fixed grid of allowed steps, so the set
$\mathcal{P}^{*}$ of prices that maximise traded volume is a short run of
neighbouring steps. The engine finds it in one pass over the submitted
limit prices and then picks one of them, $\clr = \sigma(\mathcal{P}^{*})$
\eqref{eq:clearing}, by the four-step tie-break cascade defined in
\Cref{sec:fbanotation}.

The order of those steps matters. Minimising the residual imbalance
(step 2) depends on the limit prices participants submit, so it is
exactly the lever that a late participant exploits \eqref{prop:lastmover}.
The remaining two steps --- distance to the previous clearing price
$p_{\mathrm{prev}}$, then the lowest price --- do not depend on the orders
of the current batch. Putting one of those first would remove the lever
without changing total surplus \eqref{lem:invariance}.

\subsection{Why the Price Is Uniform}

Clearing can also be seen as an optimisation problem: choose how much of each
order to execute so that total gains from trade are as large as possible,
subject to buying and selling the same quantity. The uniform price is then not
an extra rule bolted on afterwards; it falls directly out of that one
requirement, that total buying equal total selling across the whole
batch. Because the requirement applies to the batch as a whole and not
order by order, every executed order ends up trading at the same price.
Two things follow. First, the optimisation decides who trades but
not always at what price: several prices can support the same trades, and
that range is exactly $\mathcal{P}^{*}$. The price ambiguity is therefore
built into the problem, and $\sigma$ is what resolves it. Second, the
filling rule is simple. Buyers whose limit price is above the clearing price
and sellers whose limit price is below it are filled completely, orders on the
wrong side are not filled, and only orders priced exactly at the clearing
price can be partly filled. Those are the orders that must be rationed. This
holds when one price maximises the volume; when a range of prices does and the
tie-break picks a price inside it, the long side is rationed in priority order
and an order priced better than the clearing price can occasionally be filled
only in part as well (\Cref{sec:testfba}).

\subsection{Rationing Rules}
\label{sec:rationing}

At the clearing price the two sides are usually not equal, $D(\clr) \neq
S(\clr)$, so the orders on the longer side priced exactly at $\clr$ cannot all
be filled and must be rationed. The rule chosen determines what participants
do to get a bigger share. Any rule that gives more to those who increase some
variable they control encourages them to inflate that variable
\citep{haynes2020precedence}.

\begin{description}
  \item[Price-time priority within the batch.] Orders at $\clr$ are filled in
  order of submission (market orders first, then by limit price, then by
  time), which keeps continuity with continuous markets but brings back a
  speed race at the scale of the interval: if being early inside the batch
  improves the fill, participants pay to be early, which is the incentive
  the mechanism was built to remove. Under a fixed close it also pulls
  against the boundary repricing race (\Cref{sec:priceimprovement}), which
  rewards submitting late: a participant must trade queue position against
  information. This is the rule the simulated engine uses
  (\Cref{tab:fbachoices}).
  \item[Pro rata.] Each order at $\clr$ receives a share proportional to its
  size, the rule the DFBA uses after price priority. It ignores timing, but
  because a bigger order gets a bigger fill, participants submit more than they
  want to trade. This \emph{order stuffing} makes displayed sizes less
  informative.
  \item[Size priority or maximal fill.] Larger orders are filled first, or the
  number of completely filled orders is maximised. This favours large
  ``block'' orders, which may be intended, but turns pro rata's incentive to
  inflate size into a threshold effect.
  \item[Random allocation.] Fills are drawn by lottery among the orders at
  $\clr$. This removes both the timing and the size incentive and makes
  submitted size an honest statement of what the participant wants. The cost is
  an uncertainty about execution that participants cannot manage, which
  matters to anyone hedging, and the need for a random draw that participants
  can verify.
  \item[Hybrid rules.] Several of the CME's matching algorithms give part of
  the volume by time priority and the rest pro rata, balancing the two
  distortions \citep{haynes2020precedence}. A
  batch version could give a share $\kappa$ pro rata and the remaining
  $1 - \kappa$ by lottery.
\end{description}

No rule is neutral. Every rule ties the allocation to some variable, and that
variable becomes the thing participants compete on. The question is not which
rule has no distortion, but which distortion the venue prefers.

\section{Block 5: Settlement and Residual Handling}
\label{sec:settlement}

\subsection{Gross versus Net Settlement}

Gross settlement moves each matched pair separately, which is simple and
leaves a clear record of every fill. Net settlement moves only each
participant's net change in position, which reduces the number of transfers,
sometimes greatly when many participants both buy and sell in the same batch.
Where every transfer costs a fee, as on a blockchain, this is a direct saving
and one of the main practical arguments for batching
\citep{wang2018loopring,gong2023fbadex}. In a centralised venue the saving is
smaller, but reconciliation is still easier. The drawback is that netting
hides which order led to which transfer, which complicates per-fill reporting
and any fee charged per execution.

\subsection{The Residual}

The volume on the longer side that cannot be matched at $\clr$ does not
disappear, and it is not filled from any other source. What happens to it is
set by the order's time-in-force, and there are three possibilities:

\begin{enumerate}
  \item \textbf{Roll over.} The unmatched remainder moves to the next batch
  with its original limit price and, where relevant, its original submission
  time. This is the natural default for resting limit orders and is what this
  thesis implements.
  \item \textbf{Cancel.} The remainder is dropped at the end of the batch, the
  equivalent of an immediate-or-cancel order.
  \item \textbf{Cancel the whole order unless fully filled}, the equivalent of a
  fill-or-kill order. This has to be built into the clearing problem itself
  and cannot be applied afterwards.
\end{enumerate}

The third option shows where the clean structure of the clearing problem
breaks down. All-or-nothing conditions turn a problem that a simple search
over the price grid can solve into one that must test many combinations of
orders. That is why venues supporting them limit order sizes, and why
decentralised batch exchanges with rich order types use competitions between
solvers instead of one fixed algorithm \citep{harrison2025hybrid}.

\section{Summary: Design Options and Their Consequences}
\label{sec:consequences}

\Cref{tab:consequences} collects every option discussed in this chapter,
block by block, with the behaviour it is predicted to induce and the cost it
carries. The options the simulated engine implements are listed in
\Cref{tab:fbachoices}.

\begingroup
\small
\renewcommand{\arraystretch}{1.15}
\begin{longtable}{@{}>{\raggedright\arraybackslash}p{3.2cm}
  >{\raggedright\arraybackslash}p{7.2cm}
  >{\raggedright\arraybackslash}p{5.6cm}@{}}
\caption[Design options and their consequences]{Every design option of
\Cref{ch:design}, the behaviour it is predicted to induce, and its main
cost.}
\label{tab:consequences}\\
\toprule
\textbf{Option} & \textbf{Predicted consequence} & \textbf{Main cost or risk}\\
\midrule
\endfirsthead
\multicolumn{3}{@{}l}{\emph{\Cref{tab:consequences}, continued}}\\
\toprule
\textbf{Option} & \textbf{Predicted consequence} & \textbf{Main cost or risk}\\
\midrule
\endhead
\bottomrule
\endfoot
\multicolumn{3}{@{}l}{\emph{Block 1: batch interval (\Cref{sec:interval})}}\\
Short $\tau$ &
prices stay current; a speed difference $\delta$ still decides the batch with probability $\delta/\tau$ &
thin, one-sided batches that often do not cross; truthfulness weakens as orders become pivotal\\
Long $\tau$ &
speed differences become irrelevant; fuller batches &
delay of $\tau/2$ on average; stale orders; market makers cannot withdraw in time\\
Fixed close &
predictable boundary; late submission weakly dominates, producing the boundary repricing race (\Cref{sec:priceimprovement}) &
speed regains value as deadline precision; accumulation phase can empty out\\
Random close &
waiting becomes a gamble; interior optimum between information and risk of exclusion &
imprecise timing for everyone; random draw must be verifiable\\
Volume- or event-triggered close &
interval adapts to activity &
close can be triggered strategically by submitting volume\\
Random end after a minimum &
schedule predictable to within a short window; the exact last moment cannot be targeted &
a short random window gives less protection than a fully random close\\
\midrule
\multicolumn{3}{@{}l}{\emph{Block 2: information disclosure (\Cref{sec:disclosure})}}\\
Fully open book &
strong within-batch price discovery; early orders are exploited by later ones &
strongest incentive to submit late; amplifies the boundary race\\
Indicative price &
coordination on a common price estimate &
some intent leakage through price moves\\
Aggregate imbalance &
attracts offsetting liquidity without revealing limits &
weaker price anchoring\\
Sealed bid &
no exploitation of submitted orders; no informational reason to wait &
bid shading under price uncertainty, which can reduce cleared volume\\
\midrule
\multicolumn{3}{@{}l}{\emph{Block 3: order flow segregation (\Cref{sec:segregation})}}\\
Unified pool &
all liquidity meets in one auction; simplest to regulate &
uninformed flow pays for losses to informed flow\\
By participant class &
tighter prices for flow believed uninformed &
labels can be gamed; different prices for the same asset at the same time\\
By intent (DFBA) &
makers never trade with each other; no queue to pick off makers &
liquidity split across two thinner auctions; labels must be enforced\\
Maker--taker by declared intent &
simple to operate &
self-declared unless enforced\\
Maker--taker by price aggressiveness &
observable and hard to game &
penalises demand for certainty of execution\\
Maker--taker by marginality &
ties the fee to price setting &
fee unknown before the batch clears\\
Maker--taker by persistence &
rewards liquidity that stays available across batches &
needs order age tracked; every new order starts as a taker\\
\midrule
\multicolumn{3}{@{}l}{\emph{Block 4: clearing and rationing (\Cref{sec:clearing})}}\\
Tie-break led by minimum imbalance &
price responds to marginal submitted limits &
the lever for the last-mover advantage \eqref{prop:lastmover}\\
Tie-break led by a reference price &
price insensitive to marginal limits; removes the lever &
needs a reference (previous clear or external price)\\
Price-time priority &
continuity with continuous markets &
rewards being early in the batch: a speed race at the interval scale, pulling against the boundary race\\
Pro rata &
ignores timing &
order stuffing; displayed sizes less informative\\
Size priority or maximal fill &
favours block orders &
threshold incentive to inflate size\\
Random allocation &
removes timing and size incentives; size becomes an honest statement &
unmanageable execution uncertainty; verifiable randomness needed\\
Hybrid ($\kappa$ pro rata, $1-\kappa$ lottery) &
balances the two distortions &
inherits a share of both\\
\midrule
\multicolumn{3}{@{}l}{\emph{Block 5: settlement and residual (\Cref{sec:settlement})}}\\
Gross settlement &
clear record of every fill &
one transfer per matched pair\\
Net settlement &
fewer transfers; large saving where each transfer has a fee &
hides which order led to which transfer; harder per-fill fees\\
Roll over residual &
natural for resting limit orders; residual keeps its limit price &
unmatched orders accumulate across batches\\
Cancel residual (IOC) &
no stale residual &
participants must resubmit every batch\\
Fill-or-kill &
all-or-nothing execution &
turns clearing into a combinatorial problem; needs size limits or solvers\\
\end{longtable}
\endgroup

Three cross-block interactions stand out. A fixed close is costly mainly
under an open book: disclosure supplies the information that makes waiting
pay, and an imbalance-led tie-break supplies the lever to use it, so the
boundary repricing race needs Blocks~1, 2 and~4 to line up and can be
broken at any one of them. Price-time rationing and the boundary race pull
in opposite directions, early for queue position and late for
information, so neither rule removes the speed incentive on its own. And a
short $\tau$ makes every block's problems sharper at once, because thin
batches make individual orders pivotal for the price, the rationing and the
residual.
